\section{Data}
\label{sec:data}

The FAIR-MAST archive \citep{jackson2024fairmast} contains the MAST dataset of \nArchiveShots{} discharges, with each discharge stored in a separate Zarr store and accessible via HTTP. The database includes a metadata API for the discovery of available signals and shot parameters. The groups that we access from the archive are the midplane fast camera (\texttt{rbb}), the EFIT reconstruction (\texttt{efm}), the summary signals for line-averaged and Greenwald density (\texttt{esm}), the interferometer (\texttt{ane}), the core and edge Thomson scattering (\texttt{ayc}, \texttt{aye}), and the midplane $D_\alpha$ monitor (\texttt{xim}). In addition to these signals, we also access per-shot metadata including the plasma current, heating scheme, injected beam power and the interval of the flat top.

\subsection{Fast camera (\texttt{rbb})}
The \texttt{rbb} signal corresponds to the Photron SA1.1 camera (also known as the ``bullet camera B'') situated at the HM10 midplane port and aimed tangentially at the outboard edge \citep{kirk2016filaments, farley2019identification}. The recorded region is $48 \times 256$ pixels and recorded at 100{,}000 frames/s with 10 bits resolution. The offset of this region on the sensor is stored as part of the per-shot metadata. The exposure times vary between 1 and 10~$\mu$s, and the most common settings are 4, 7, and 10~$\mu$s. Also recorded are the lens, aperture, and filter settings and whether the camera was inverted. We do not rely on the orientation of the camera for our detection pipeline (Section~\ref{sec:method}), but all camera configurations are catalogued in Appendix~\ref{app:views}, along with the details of their calibrations.

\subsection{Discharge selection}
\nStripShots{} discharges from FAIR-MAST have the fast camera recording a $48 \times 256$ pixel region. These are in the form of bursts, high-speed recordings made in a window within a slower recording of the whole discharge at 1--2~kHz. The median duration of these bursts is \burstMedianMs{}~ms, with extremes of \burstMinMs{} to \burstMaxMs{}~ms. Of these, \nStripTenUs{} have frame intervals of 10~$\mu$s. We select discharges for further processing with the following criteria: (i) the burst of interest overlaps the plasma flat top, extended by 20~ms either side, by at least 2{,}000 fast frames. We find \nOkBurst{} of these discharges. (ii) the flat top must be defined and at least 50~ms long. We find \nOkFlatTop. (iii) a peak plasma current of at least \ipCandMin{}~kA. We find \nOkCurrent{} discharges, and \nShotsCandidates{} discharges satisfy all three conditions.

For the \nShotsFetched{} selected discharges, we download and process the frames that are in the burst and between the extended flat top. The median length of this processing window is \windowMedianMs{}~ms, leading to a total of \nFramesProcessedM{} million frames processed. Sixteen of the selected discharges have a slower burst rate of \slowFpsKhz{}~kHz; these are included in the processing, but are excluded from the statistics to avoid changing time resolution. The set of processed discharges is shown broken down by campaign in Table~\ref{tab:campaigns}, with the number of discharges, exposure settings, pixel scales, and median velocity and width for each.

\begin{table}[t]
\centering
\scriptsize
\setlength{\tabcolsep}{4pt}\begin{tabular}{llrrllllr}
\toprule
Campaign & Years & Shots & In stats & Exp. ($\mu$s) & Pooled scale (mm/px) & Own scales & $\tilde v_r$ (m/s) & $\tilde\delta_r$ (cm) \\
\midrule
M6 & 2006--2007 & 257 & 220 & 1, 5, 7, 10 & 3.7 (3.3--4.1) & 3 (2.9--3.7) & 419 & 3.6 \\
M7 & 2008--2009 & 15 & 13 & 4, 10 & 3.3 (3.0--3.8) & 1 (3.3) & 451 & 3.6 \\
M8 & 2011 & 32 & 25 & 2, 4 & 5.2 (3.6--11.5) & 1 (5.2) & 627 & 5.2 \\
M9 & 2013 & 48 & 36 & 4, 10 & 2.9 (2.5--4.2) & 1 (2.7) & 323 & 2.8 \\
\bottomrule
\end{tabular}

\caption{Processed discharges, broken down by campaign, and within these by exposure setting. The ``In stats'' column shows only those discharges with at least 50 selected tracks. The ``pooled'' and ``own'' columns are the calibrations' scales presented in Section~\ref{sec:calibration}, with their 95\% bootstrap intervals. The columns $\tilde v_r$ and $\tilde\delta_r$ are the median over discharges of each discharge's per-discharge median velocity and width.}
\label{tab:campaigns}
\end{table}

The LCFS boundary contour, $q_{95}$, plasma current, magnetic axis, and vacuum toroidal field at the magnetic axis are taken from the EFIT reconstruction, sampled at 5~ms. The LCFS radius in the outboard midplane is taken as the maximum major radius that is within 5~cm of the magnetic axis height. This is the quantity that the camera limb tracks. The line averaged density is taken from the \texttt{esm} summary signal when available, otherwise it is taken from the CO$_2$ interferometer summary. The Greenwald fraction uses the Greenwald density $n_{\rm GW} = I_p/(\pi a^2)$, with plasma current $I_p$ in MA and minor radius $a$ from EFIT. Thomson scattering profiles are only available in the later campaigns and only reach the LCFS in a subset of the discharges within the flat top, \nTeMeasured{} of the \nShotsStats{} discharges used for statistical summaries.
